%%%%%%%%%%%%%%%%%%%%%%%%%%%%%%%%%%%%%%%%%%%%%%%%%%%%%%%%%%%%%%%%%%%%%%%%%%%%%%%%
% ECE M237: Dynamic Programming -- Course Project
% Topic Introduction & Literature Review
% Hybrid Model-Based and Data-Driven Score Fusion for Robot Trajectory Generation
%%%%%%%%%%%%%%%%%%%%%%%%%%%%%%%%%%%%%%%%%%%%%%%%%%%%%%%%%%%%%%%%%%%%%%%%%%%%%%%%

\documentclass[letterpaper, 10pt, conference]{ieeeconf}

\IEEEoverridecommandlockouts
\overrideIEEEmargins

\usepackage{times}
\usepackage{color, colortbl, xcolor}

\let\proof\relax
\let\endproof\relax
\let\example\relax
\let\endexample\relax

\usepackage{amsmath}
\usepackage{amsthm}
\usepackage{cases}
\usepackage{amsfonts}
\usepackage{amssymb}
\usepackage{mathtools}
\usepackage{mathrsfs}
\usepackage{nicefrac}

\usepackage{multicol}
\usepackage{multirow}
\usepackage[bookmarks=true]{hyperref}
\usepackage{cleveref}
\usepackage[pdftex]{graphicx}
\usepackage{cancel}
\let\labelindent\relax
\usepackage{enumitem}
\usepackage{comment}
\usepackage{booktabs}
\usepackage{array}

\newtheorem{theorem}{Theorem}
\newtheorem{proposition}{Proposition}
\newtheorem{definition}{Definition}
\newtheorem{remark}{Remark}

% Convenience macros
\newcommand{\R}{\mathbb{R}}
\newcommand{\E}{\mathbb{E}}
\newcommand{\norm}[1]{\left\|#1\right\|}
\newcommand{\KL}{\mathrm{KL}}
\newcommand{\dd}{\mathrm{d}}
\newcommand{\score}{\nabla_{\tau_t}\log p_t}
\newcommand{\scoredata}{s_\theta(\tau_t,t)}
\newcommand{\scoremodel}{\nabla_{\tau_t}\log p_t^{\mathrm{model}}(\tau_t)}

\title{\LARGE \bf
Hybrid Model-Based and Data-Driven Score Fusion\\
for Robot Trajectory Generation}

\author{%
  \parbox{3in}{\centering Zihao Sun\\
               UCLA ECE, UID: 706-467-647}
  \thanks{ECE M237: Dynamic Programming, Spring 2026 --- Topic Introduction \& Literature Review.}
}

\begin{document}

\maketitle

%%%%%%%%%%%%%%%%%%%%%%%%%%%%%%%%%%%%%%%%%%%%%%%%%%%%%%%%%%%%%%%%%%%%%%%%%%%%%%%%
\begin{abstract}
We investigate the fusion of two complementary sources of information for robot
trajectory generation: (1) a \emph{data-driven score} learned from expert
demonstrations via flow matching, and (2) a \emph{model-based score} computed
analytically from the robot's dynamics and cost function.  We formalize this
combination as a \emph{product-of-experts} decomposition and study the resulting
trajectory distribution.  In the linear-quadratic (LQ) special case, we conjecture
that the model-based correction reduces exactly to the LQR Riccati gain---an
elegant bridge between generative modeling and classical dynamic programming.
This document provides the theoretical background, literature survey, and a
roadmap for our project.
\end{abstract}

%%%%%%%%%%%%%%%%%%%%%%%%%%%%%%%%%%%%%%%%%%%%%%%%%%%%%%%%%%%%%%%%%%%%%%%%%%%%%%%%
\section{The Big Picture}
\label{sec:bigpicture}

In robot motion planning, two complementary information sources are often available:
\begin{itemize}[leftmargin=*]
  \item \textbf{Data:} a dataset of expert demonstrations encoding diverse,
        multimodal behaviors.
  \item \textbf{Physics:} a dynamics model $x_{k+1} = f(x_k, u_k)$ together
        with a cost $J(\tau)$ encoding the task objective.
\end{itemize}
Recent generative models---diffusion models and flow matching---have shown
impressive results learning trajectory distributions from data~\cite{janner2022,chi2023diffusionpolicy}.
Separately, Model-Based Diffusion (MBD)~\cite{mbd} computes the score
function analytically from the dynamics model, without any data.  These two
paradigms have remained disjoint.

\textbf{Our project:} can we \emph{fuse} the data-driven score (capturing
multimodal demonstrations) with the model-based score (guaranteeing feasibility
and optimality) to obtain the best of both?  Does this fusion connect to
classical optimal control---specifically, does the model-based correction
reduce to the LQR Riccati gain in the LQ case?  This is centrally a question
about \textbf{merging offline and online methods}---the theme of ECE M237.

%%%%%%%%%%%%%%%%%%%%%%%%%%%%%%%%%%%%%%%%%%%%%%%%%%%%%%%%%%%%%%%%%%%%%%%%%%%%%%%%
\section{Theoretical Background}
\label{sec:theory}

This section develops the mathematical foundations underlying our project:
(i) optimal control theory and dynamic programming,
(ii) score-based and flow matching generative models, and
(iii) reinforcement learning and approximate dynamic programming.

%---------------------------------------------------------------------
\subsection{Optimal Control and Dynamic Programming}
\label{sec:oc}

\subsubsection{Problem Formulation}
Consider a discrete-time dynamical system
\begin{equation}
  x_{k+1} = f(x_k, u_k), \quad x_0 \in \R^n,\; u_k \in \mathcal{U} \subseteq \R^m,
  \label{eq:dynamics}
\end{equation}
and a finite-horizon cost functional
\begin{equation}
  J(\pi) = \E\!\left[\sum_{k=0}^{H-1} g(x_k, u_k) + V_H(x_H)\right],
  \label{eq:cost}
\end{equation}
where $g:\R^n\times\R^m\to\R$ is the stage cost and $V_H$ is the terminal cost.
The goal is to find a \emph{policy} $\pi = \{\mu_0,\ldots,\mu_{H-1}\}$, with
$u_k = \mu_k(x_k)$, that minimizes $J(\pi)$.

\subsubsection{Principle of Optimality and Bellman Equation}
Bellman's Principle of Optimality states that every sub-policy of an optimal
policy is itself optimal with respect to the sub-problem originating at that
point~\cite{bertsekas2012dp}.  This gives rise to the \emph{Bellman equation}
(backward recursion):
\begin{equation}
  V_k(x) = \min_{u \in \mathcal{U}}\Bigl[g(x,u) + V_{k+1}(f(x,u))\Bigr],\quad
  V_H(x) \text{ given},
  \label{eq:bellman}
\end{equation}
where $V_k:\R^n\to\R$ is the \emph{value function} (optimal cost-to-go) at
stage $k$.  The optimal policy is recovered as
$\mu_k^*(x) = \arg\min_{u}[g(x,u) + V_{k+1}(f(x,u))]$.

\subsubsection{Linear Quadratic Regulator (LQR)}
For linear dynamics $x_{k+1} = Ax_k + Bu_k$ and quadratic cost
$g(x,u) = x^\top Q x + u^\top R u$ (with $Q \succeq 0$, $R \succ 0$,
$V_H(x) = x^\top Q_H x$), the Bellman equation admits a closed-form solution.
Assuming $V_k(x) = x^\top P_k x$, substituting into
\eqref{eq:bellman} and minimizing over $u$ gives the
\emph{discrete-time Riccati difference equation} (DARE):
\begin{align}
  P_k &= A^\top P_{k+1} A + Q
         - A^\top P_{k+1} B\bigl(R + B^\top P_{k+1} B\bigr)^{-1}
           B^\top P_{k+1} A,
  \label{eq:riccati}\\
  K_k &= \bigl(R + B^\top P_{k+1} B\bigr)^{-1}B^\top P_{k+1} A,
  \label{eq:gain}
\end{align}
with terminal condition $P_H = Q_H$.  The optimal control law is the linear
state-feedback $u_k^* = -K_k x_k$.  For infinite-horizon problems with stable
$(A,B)$, the sequence $P_k$ converges to the unique positive definite solution
$P_\infty$ of the \emph{algebraic Riccati equation}
$P = A^\top P A + Q - A^\top P B(R+B^\top P B)^{-1}B^\top P A$.

\begin{remark}[Connection to Score Functions]
  A key conjecture of our project is that, for LQ systems, the gradient of the
  model-based log-density satisfies
  $\scoremodel = -P_k x$,
  where $P_k$ is the Riccati solution \eqref{eq:riccati}.  This would mean the
  model-based score correction in the generative process \emph{is} the LQR
  feedback gain.
\end{remark}

\subsubsection{Hamilton--Jacobi--Bellman Equation (Continuous Time)}
In continuous time, let $\dot{x} = f(x,u)$ with cost
$J = \int_0^T g(x,u)\,\dd t + V_T(x(T))$.  The value function
$V(x,t) = \min_u J$ satisfies the \emph{HJB PDE}:
\begin{equation}
  -\frac{\partial V}{\partial t}(x,t)
  = \min_{u\in\mathcal{U}}\Bigl[g(x,u) + \nabla_x V(x,t)^\top f(x,u)\Bigr],
  \label{eq:hjb}
\end{equation}
with terminal condition $V(x,T) = V_T(x)$.  For stochastic dynamics
$\dd x = f(x,u)\,\dd t + \sigma\,\dd W_t$, the HJB equation acquires a
diffusion term $\frac{\sigma^2}{2}\Delta_x V$, where $\Delta_x$ is the
Laplacian.  Solving the HJB equation is subject to the \emph{curse of
dimensionality}: the state space grows exponentially in dimension, motivating
the approximate methods discussed in Section~\ref{sec:rl}.

\subsubsection{Contraction Mapping and Value Iteration}
The Bellman operator $\mathcal{T}[V](x) = \min_u[g(x,u) + V(f(x,u))]$ is a
$\gamma$-contraction in the $\sup$-norm (for discounted problems with discount
factor $\gamma < 1$) by Banach's Fixed-Point Theorem.  \emph{Value iteration}
$V_{k+1} \leftarrow \mathcal{T}[V_k]$ therefore converges geometrically to
the unique fixed point $V^*$.  \emph{Policy iteration}---alternating between
policy evaluation (solving the Bellman equation for a fixed policy) and policy
improvement (one-step greedy update)---converges in a finite number of steps
for finite MDPs and achieves superlinear convergence~\cite{bertsekas2012dp}.

%---------------------------------------------------------------------
\subsection{Score-Based Generative Models and Flow Matching}
\label{sec:fm}

\subsubsection{Score Functions and Diffusion Models}
Let $p_0$ be a data distribution over trajectories $\tau\in\R^d$.  A diffusion
model defines a \emph{forward SDE} that gradually corrupts $\tau$ toward
Gaussian noise:
\begin{equation}
  \dd\tau_t = f_t(\tau_t)\,\dd t + g_t\,\dd W_t,\quad t\in[0,T],
  \label{eq:forward_sde}
\end{equation}
and a \emph{reverse SDE} (Anderson 1982) that generates samples by denoising:
\begin{equation}
  \dd\tau_t = \bigl[f_t(\tau_t) - g_t^2\score(\tau_t)\bigr]\dd t
              + g_t\,\dd\bar{W}_t,
  \label{eq:reverse_sde}
\end{equation}
where $\score = \nabla_{\tau_t}\log p_t(\tau_t)$ is the \emph{score function}
of the marginal $p_t$~\cite{song2021score}.  In standard diffusion models, a
neural network $s_\theta(\tau_t,t)$ approximates the score via
\emph{denoising score matching}~\cite{ho2020ddpm}:
\begin{equation}
  \mathcal{L}_{\mathrm{DSM}}(\theta)
  = \E_{t,\tau_0,\tau_t}\!\bigl[\norm{s_\theta(\tau_t,t)
    - \nabla_{\tau_t}\log p_{t|0}(\tau_t|\tau_0)}^2\bigr].
  \label{eq:dsm}
\end{equation}
Classifier guidance~\cite{dhariwal2021diffusion} augments the score with the
gradient of a cost function $c(\tau_t)$:
$\tilde{s}(\tau_t,t) = s_\theta(\tau_t,t) + \lambda\nabla_{\tau_t}\log c(\tau_t),$
enabling conditional generation without retraining.

\subsubsection{Flow Matching}
Flow Matching (FM)~\cite{lipman2023flow} replaces the stochastic SDE with a
deterministic ODE
\begin{equation}
  \dot{x}_t = v_\theta(x_t,t),\quad x_0\sim\mathcal{N}(0,I),\; t\in[0,1].
  \label{eq:fm_ode}
\end{equation}
The \emph{conditional flow matching} (CFM) objective regresses the velocity
field on simple conditional paths (typically linear interpolants
$x_t = (1-t)x_0 + tx_1$):
\begin{equation}
  \mathcal{L}_{\mathrm{CFM}}(\theta)
  = \E_{t,(x_0,x_1)}\!\bigl[\norm{v_\theta(x_t,t) - (x_1-x_0)}^2\bigr].
  \label{eq:cfm}
\end{equation}
FM has several advantages over diffusion for robotics: deterministic
generation (same noise $\Rightarrow$ same trajectory), faster inference (single
ODE integration), and compatibility with optimal control theory since the ODE
has a clean variational interpretation.

\subsubsection{Guided Flow Matching via Optimal Control}
\textit{OC-Flow}~\cite{ocflow} casts guided generation as an optimal
control problem: starting from a pre-trained FM prior $v_\theta$, find a
\emph{control signal} $u_t$ that minimizes a terminal reward $R(x_1)$ while
penalizing deviation from the prior:
\begin{equation}
  \min_{u_t} \;-R(x_1) + \tfrac{1}{2}\int_0^1\norm{u_t}^2\,\dd t,\quad
  \dot{x}_t = v_\theta(x_t,t) + u_t.
  \label{eq:ocflow}
\end{equation}
The necessary optimality conditions (Pontryagin Maximum Principle) give the
co-state $\lambda_t$ satisfying $\dot{\lambda}_t = -(\partial v_\theta/\partial x_t)^\top\lambda_t$,
and the optimal control $u_t^* = \lambda_t$.  This framework provides a
principled, training-free way to guide flow matching with an arbitrary
differentiable cost.

\subsubsection{Stochastic Optimal Control and Score Matching}
SOCM~\cite{domingoenrich2024socm} reveals a deep connection: the optimal
\emph{stochastic} control that steers the SDE \eqref{eq:forward_sde} from a
prior to a target distribution is intimately related to score matching.
Specifically, the optimal control $u^*$ can be learned by minimizing a
least-squares objective analogous to the denoising score matching loss
\eqref{eq:dsm}, with the path-wise reparameterization trick reducing variance.
This unifies score-based generative modeling and stochastic optimal control
under a single mathematical umbrella.

%---------------------------------------------------------------------
\subsection{Reinforcement Learning and Approximate Dynamic Programming}
\label{sec:rl}

\subsubsection{Markov Decision Processes}
A \emph{Markov Decision Process} (MDP) is a tuple
$(\mathcal{S},\mathcal{A},P,r,\gamma)$ where $\mathcal{S}$ is the state space,
$\mathcal{A}$ the action space, $P(s'|s,a)$ the transition kernel,
$r(s,a)$ the reward function, and $\gamma\in(0,1)$ the discount factor.
The optimal \emph{action-value function} $Q^*(s,a)$ satisfies the
\emph{Bellman optimality equation}:
\begin{equation}
  Q^*(s,a) = r(s,a) + \gamma\sum_{s'}P(s'|s,a)\max_{a'}Q^*(s',a').
  \label{eq:bellman_q}
\end{equation}
For continuous state spaces, \eqref{eq:bellman_q} cannot be solved exactly,
motivating function approximation.

\subsubsection{Value Function Approximation and Fitted Methods}
\emph{Linear function approximation} represents $V^\pi(s)\approx\phi(s)^\top w$
with feature vector $\phi(s)$.  \emph{Least-squares policy evaluation} (LSPE)
solves the projected Bellman equation exactly in the feature subspace.
\emph{Fitted value iteration} replaces the tabular update with supervised
regression: given a dataset $\{(s_i,a_i,r_i,s_i')\}$, fit
$V_{\theta}\leftarrow\arg\min_\theta\sum_i(V_\theta(s_i) - [r_i + \gamma\max_{a'}V_\theta(s_i')])^2$.
Similarly, \emph{fitted Q-iteration} (FQI) fits the Q-function.

\subsubsection{Deep Reinforcement Learning}
\textbf{DQN}~\cite{mnih2015dqn} parametrizes $Q^*(s,a)$ with a deep neural
network and stabilizes training via (i) a \emph{replay buffer} to break temporal
correlations and (ii) a \emph{target network} that provides stable regression
targets.  This enabled superhuman performance on Atari games directly from pixels.

\textbf{Actor-Critic Methods} maintain both a \emph{critic} $V_\theta(s)$ (or
$Q_\theta(s,a)$) and an \emph{actor} $\pi_\phi(a|s)$.  The actor is updated
via policy gradient:
$\nabla_\phi J(\phi) = \E_{\pi_\phi}[Q_\theta(s,a)\nabla_\phi\log\pi_\phi(a|s)].$
\textbf{Soft Actor-Critic (SAC)}~\cite{haarnoja2018sac} augments the reward
with an entropy bonus $\alpha\mathcal{H}(\pi(\cdot|s))$, yielding a maximum
entropy policy that is robust to model errors and hyperparameter choices.

\textbf{Policy Gradient / PPO}~\cite{schulman2017ppo} directly optimizes
$J(\phi)=\E_\pi[\sum_t r_t]$ by gradient ascent.  Proximal Policy
Optimization (PPO) constrains policy updates via a clipped surrogate
objective, providing sample efficiency and stability.

\subsubsection{Exploration and Exploitation}
The exploration--exploitation dilemma is fundamental in RL~\cite{suttbarto2018rl}.
In our project, the blending weight $\lambda$ plays an analogous role:
$\lambda=0$ exploits the data prior (offline, no exploration); $\lambda\to\infty$
explores via the model (online, potentially unimodal).  This connection is made
precise in Section~\ref{sec:proposal}.

\subsubsection{Online Search: Rollout and MPC}
\emph{Rollout} improves a base policy $\mu$ by one-step look-ahead: at each
state, evaluate several candidate actions by simulating $\mu$ forward and
take the best.  \emph{Model Predictive Control} (MPC) solves a
finite-horizon optimization at each step and executes only the first action
(receding horizon).  MBD~\cite{mbd} can be interpreted as
sampling-based MPC in score space, and adding the data prior corresponds to
warm-starting MPC with an offline policy.

%%%%%%%%%%%%%%%%%%%%%%%%%%%%%%%%%%%%%%%%%%%%%%%%%%%%%%%%%%%%%%%%%%%%%%%%%%%%%%%%
\section{Background: Diffusion and Flow Matching for Trajectories}
\label{sec:diffusion_background}

%---------------------------------------------------------------------
\subsection{Diffusion Models for Trajectory Planning}

A diffusion model for trajectory planning treats each expert trajectory
$\tau = (q_0,q_1,\ldots,q_H)$ as a high-dimensional data point.  The forward
process \eqref{eq:forward_sde} corrupts $\tau$ over time $t\in[0,T]$ toward a
Gaussian prior.  The reverse process \eqref{eq:reverse_sde} generates new
trajectories by starting from Gaussian noise and iteratively denoising with
the learned score $s_\theta(\tau_t,t)$.

\textbf{Diffuser}~\cite{janner2022}: the foundational work. Models the
joint state-action trajectory as a high-dimensional ``image'' using a temporal
U-Net, and uses classifier guidance for goal-conditioning. Demonstrated on D4RL
offline RL benchmarks (maze navigation, locomotion). No dynamics awareness.

\textbf{Diffusion Policy}~\cite{chi2023diffusionpolicy}: applies diffusion to
action-space policy learning for manipulation. FiLM conditioning on observations.
Limitation: action-space diffusion lacks trajectory-level feasibility guarantees.

\textbf{MPD}~\cite{carvalho2023mpd}: learns a trajectory prior, then samples
the posterior conditioned on cost guidance during denoising. B-spline
parameterization for smoothness. Closest data-driven precedent to our work, but
does not use the dynamics model for score computation.

%---------------------------------------------------------------------
\subsection{Flow Matching for Robot Motion Planning}

\textbf{FlowMP}~\cite{flowmp}: extends FM to second-order trajectory
dynamics (acceleration), producing smooth, physically executable motions.
Achieves $2\times$ faster inference than MPD. Purely data-driven.

\textbf{SafeFlow}~\cite{safeflow}: integrates FM with \emph{flow
matching barrier functions} (FMBF), enabling training-free safety enforcement
at test time. 100\% safety rate on 2D navigation and 7-DoF Panda manipulation.
Limitation: safety $\ne$ dynamics feasibility.

\textbf{SafeFlowMatcher}~\cite{yang2025safeflowmatcher}: a two-phase
prediction-correction integrator. Prediction integrates the FM flow; correction
refines with a CBF-based QP. Achieves faster, smoother, safer paths than diffusion baselines.

\textbf{UniConFlow}~\cite{yang2025uniconflow}: unified constrained FM enforcing
both equality and inequality constraints via prescribed-time zeroing functions
and QP-based guidance. State of the art on certified motion planning metrics.

\textbf{GoalFlow}~\cite{xing2025goalflow}: goal-driven FM for autonomous
driving. Single denoising step, state-of-the-art on the NAVSIM benchmark.

%---------------------------------------------------------------------
\subsection{The Limitation: No Dynamics Awareness}

All methods above learn the score or velocity field \emph{entirely from data}.
Generated trajectories can violate dynamics constraints, degrade under
distribution shift, and have no formal connection to the optimal control
solution. This gap is exactly what MBD and our project address.

%%%%%%%%%%%%%%%%%%%%%%%%%%%%%%%%%%%%%%%%%%%%%%%%%%%%%%%%%%%%%%%%%%%%%%%%%%%%%%%%
\section{Model-Based Diffusion (MBD)}
\label{sec:mbd}

%---------------------------------------------------------------------
\subsection{Core Idea}

MBD~\cite{mbd} takes the opposite approach: it \emph{computes} the score
analytically from the dynamics model instead of learning it.  Given a trajectory
optimization problem $\min_\tau J(\tau)$ subject to $q_{k+1} = f(q_k,u_k)$, MBD
defines the target distribution
\begin{equation}
  p_0(\tau) \propto \exp\!\bigl(-J(\tau)/\lambda\bigr)\cdot\mathbf{1}[\tau\text{ is dynamically feasible}],
  \label{eq:mbd_target}
\end{equation}
and estimates the score $\nabla_\tau\log p_0(\tau)$ via \emph{importance-weighted
Monte Carlo rollouts}: sample many perturbations of the current trajectory,
roll them through the dynamics model, weight by $\exp(-J(\tau)/\lambda)$, and
average the directions. This is sampling-based trajectory optimization in score space.

%---------------------------------------------------------------------
\subsection{Strengths and Limitations}

\textbf{Strengths:} No training data required; generalizes to new tasks and
robots without retraining; guarantees dynamic feasibility by construction;
outperforms PPO by 59\% on MuJoCo locomotion.

\textbf{Limitations:} Cannot capture multimodal trajectory distributions (no
data prior); computationally expensive (many MC rollouts); the paper
(Section~4.3) mentions that data ``can be naturally integrated'' but
\emph{does not formalize this integration}---no product-of-experts
decomposition, no convergence analysis, no tradeoff bound.  This informal
mention is \textbf{exactly the gap our project fills}.

%%%%%%%%%%%%%%%%%%%%%%%%%%%%%%%%%%%%%%%%%%%%%%%%%%%%%%%%%%%%%%%%%%%%%%%%%%%%%%%%
\section{The Score Function and Its Two Sources}
\label{sec:scores}

%---------------------------------------------------------------------
\subsection{Data-Driven Score}

Learned from demonstrations via a neural network:
$\scoredata \approx \nabla_{\tau_t}\log p_t^{\mathrm{data}}(\tau_t).$
\textit{Knows:} multimodal expert behavior, implicit constraints.
\textit{Does not know:} robot dynamics, new-environment costs.

%---------------------------------------------------------------------
\subsection{Model-Based Score}

Computed from dynamics and cost:
$\scoremodel.$
\textit{Knows:} physics, feasibility, optimality.
\textit{Does not know:} demonstration distribution, multimodal behaviors.

%---------------------------------------------------------------------
\subsection{Product-of-Experts Fusion}

Can we combine them?
\begin{equation}
  \nabla_{\tau_t}\log p_t^{\mathrm{target}}
  = \underbrace{\scoredata}_{\text{offline, from data}}
  + \lambda\underbrace{\scoremodel}_{\text{online, from dynamics}}.
  \label{eq:poe_score}
\end{equation}
This is a \emph{product-of-experts} decomposition:
\begin{equation}
  p_t^{\mathrm{target}}(\tau_t)
  \propto p_t^{\mathrm{data}}(\tau_t)\cdot
  \bigl[p_t^{\mathrm{model}}(\tau_t)\bigr]^\lambda.
  \label{eq:poe}
\end{equation}
The weight $\lambda$ controls the offline--online tradeoff:
$\lambda=0$ (pure data), $\lambda=1$ (equal weight), $\lambda\to\infty$ (pure model).
This is mathematically analogous to classifier guidance~\cite{dhariwal2021diffusion}
and CFGRL~\cite{cfgrl}, but with an analytically computed
model-based score instead of a learned classifier.

%---------------------------------------------------------------------
\subsection{Behavioral Score Diffusion (Opposite Direction)}

BSD~\cite{bsd} goes the \emph{opposite} direction from MBD: it estimates
scores entirely from data via kernel-weighted Nadaraya-Watson estimation,
with no dynamics model. Triple-kernel weighting (diffusion proximity, state
context, goal relevance). Achieves 98.5\% of MBD's reward without a model.
\textbf{BSD confirms the field has split into model-only vs.\ data-only camps;
nobody occupies the middle ground.}

%%%%%%%%%%%%%%%%%%%%%%%%%%%%%%%%%%%%%%%%%%%%%%%%%%%%%%%%%%%%%%%%%%%%%%%%%%%%%%%%
\section{The Gap}
\label{sec:gap}

Table~\ref{tab:gap} summarizes the landscape.

\begin{table}[h]
  \centering
  \caption{Landscape of trajectory generation methods by score source.}
  \label{tab:gap}
  \small
  \begin{tabular}{lccc}
    \toprule
    \textbf{Method} & \textbf{Year} & \textbf{Data?} & \textbf{Model?} \\
    \midrule
    Diffuser~\cite{janner2022}            & 2022 & \checkmark & $\times$ \\
    Diffusion Policy~\cite{chi2023diffusionpolicy}& 2023 & \checkmark & $\times$ \\
    MPD~\cite{carvalho2023mpd}                    & 2024 & \checkmark & partial \\
    MBD~\cite{mbd}                         & 2024 & optional*  & \checkmark \\
    FlowMP~\cite{flowmp}                & 2025 & \checkmark & $\times$ \\
    SafeFlow~\cite{safeflow}               & 2025 & \checkmark & safety only \\
    OC-Flow~\cite{ocflow}                 & 2025 & \checkmark & reward only \\
    S$^2$Diff~\cite{s2diff}              & 2025 & $\times$   & \checkmark \\
    BSD~\cite{bsd}                          & 2026 & \checkmark & $\times$ \\
    \textbf{Ours (proposed)}                      & 2026 & \checkmark & \checkmark \\
    \bottomrule
  \end{tabular}
  \\[4pt]
  \footnotesize *MBD mentions data integration informally; does not formalize it.
\end{table}

The gap is clear: nobody has formalized the product-of-experts fusion, studied
the $\lambda$-tradeoff, or connected the model-based correction to the Riccati
equation.

%%%%%%%%%%%%%%%%%%%%%%%%%%%%%%%%%%%%%%%%%%%%%%%%%%%%%%%%%%%%%%%%%%%%%%%%%%%%%%%%
\section{Our Proposal: Product-of-Experts Score Fusion}
\label{sec:proposal}

%---------------------------------------------------------------------
\subsection{The Algorithm}

At each step of the reverse flow ODE, we use the composite score
\eqref{eq:poe_score}.  $s_\theta$ is the learned flow matching score
from the offline-trained model; $\scoremodel$ is computed online via Monte Carlo
rollouts (nonlinear systems) or analytically (LQ systems); $\lambda\geq 0$
is the blending weight.  Optionally, safety is incorporated as a third term:
\begin{equation}
  \nabla_{\tau_t}\log p_t^{\mathrm{target}}
  = s_\theta + \lambda_{\mathrm{model}}\cdot\scoremodel
    + \lambda_{\mathrm{safe}}\cdot\nabla\log p_t^{\mathrm{safe}},
  \label{eq:three_score}
\end{equation}
where $\nabla\log p_t^{\mathrm{safe}}$ encodes collision avoidance via CBFs,
connecting to SafeFlow~\cite{safeflow} and
S$^2$Diff~\cite{s2diff}.

%---------------------------------------------------------------------
\subsection{The LQ Special Case: Conjectured Riccati Connection}

For linear dynamics $x_{k+1} = Ax_k + Bu_k$ and quadratic cost
$J = \sum_k(x_k^\top Qx_k + u_k^\top Ru_k)$, we conjecture:
\begin{equation}
  \scoremodel = -P_k x_t,
  \label{eq:riccati_score}
\end{equation}
where $P_k$ satisfies the DARE \eqref{eq:riccati}.  If this holds,
the model-based correction in the generative model is \emph{exactly the LQR
optimal feedback gain}---the object derived via dynamic programming in the
course.  This is a surprising and elegant bridge between generative modeling
and classical control.

%---------------------------------------------------------------------
\subsection{Plan of Work}

\textbf{Theory:} (1) Verify the LQ-Riccati connection analytically.
(2) If tractable, prove convergence of the hybrid ODE in the LQ case.
(3) Characterize the $\lambda$-tradeoff: how does blending affect cost vs.
feasibility?

\textbf{Implementation:} (1) Train a flow matching model on expert trajectories
(2D navigation + manipulator). (2) Implement the model-based score (analytical
for LQ, Monte Carlo for nonlinear). (3) Implement the hybrid algorithm with
tunable $\lambda$. (4) Run experiments: feasibility comparison, $\lambda$-sweep
tradeoff curve, robustness to distribution shift, LQR ablation.

\textbf{Fallback:} empirical $\lambda$-tradeoff characterization + LQ-Riccati
connection as the main theoretical observation.

%%%%%%%%%%%%%%%%%%%%%%%%%%%%%%%%%%%%%%%%%%%%%%%%%%%%%%%%%%%%%%%%%%%%%%%%%%%%%%%%
\section{Connection to Dynamic Programming and Course Material}
\label{sec:dp_connection}

\begin{table}[h]
  \centering
  \caption{Connections to ECE M237 lecture topics.}
  \label{tab:dp}
  \small
  \begin{tabular}{p{0.8cm}p{3.0cm}p{3.5cm}}
    \toprule
    \textbf{Lec} & \textbf{Topic} & \textbf{Connection} \\
    \midrule
    3 & System models, Lyapunov stability
      & Dynamics model defines $p^{\mathrm{model}}$; Lyapunov certifies stability. \\
    4 & Bellman, value/policy iteration, curse of dimensionality
      & Model-based score $=$ gradient of value function; curse of dimensionality motivates learned prior. \\
    5 & LQR, Riccati, HJB PDE
      & \textbf{Deepest connection.} Model-based score $= -P_k x$ (Riccati gain); HJB governs continuous-time model score. \\
    6 & Exploration-exploitation
      & $\lambda$ is the exploration-exploitation tradeoff weight. \\
    7--8 & Approx.\ DP, deep RL
      & FM model = function approximator for trajectory distribution; fitted value iteration $\approx$ training FM on rollout data. \\
    9--10 & MPC, rollout, MCTS
      & MBD MC score estimation $=$ sampling-based MPC in score space; data prior $=$ warm-starting MPC. \\
    11--12 & Safety, reachability, CBFs
      & Safety enters as third score term $\nabla\log p^{\mathrm{safe}}$. \\
    \bottomrule
  \end{tabular}
\end{table}

%%%%%%%%%%%%%%%%%%%%%%%%%%%%%%%%%%%%%%%%%%%%%%%%%%%%%%%%%%%%%%%%%%%%%%%%%%%%%%%%
\section{Literature Review}
\label{sec:litreview}

%---------------------------------------------------------------------
\subsection{Foundational Works: Diffusion for Planning}

\textbf{DDPM}~\cite{ho2020ddpm} introduced the denoising diffusion probabilistic
model framework, connecting diffusion probabilistic models to denoising score
matching with Langevin dynamics.  This is the theoretical bedrock for all
diffusion-based planners.

\textbf{Score SDE}~\cite{song2021score} unified all noise perturbation processes
under a single SDE framework, providing continuous-time score matching and
enabling deterministic (ODE-based) sampling.  This formulation is critical for
connecting diffusion models to optimal control.

\textbf{Diffuser}~\cite{janner2022}: first to apply diffusion to
trajectory planning. Temporal U-Net, classifier guidance for goal-conditioning.
D4RL benchmarks. No dynamics awareness.

\textbf{Diffusion Policy}~\cite{chi2023diffusionpolicy}: action-space diffusion
for manipulation; spawned a large ecosystem. Lacks trajectory-level feasibility.

\textbf{MPD}~\cite{carvalho2023mpd}: MAP formulation with cost guidance during
denoising, B-spline smoothness. Closest data-driven precedent; no dynamics-based
score.

%---------------------------------------------------------------------
\subsection{Model-Based Score Computation}

\textbf{MBD}~\cite{mbd}: computes score analytically from dynamics via
Monte Carlo importance sampling. Target distribution \eqref{eq:mbd_target}.
Outperforms PPO by 59\% on MuJoCo; mentions data integration informally.
\textbf{The gap our project fills.}

\textbf{BSD}~\cite{bsd}: kernel-weighted Nadaraya-Watson score estimation
from data only; no dynamics model. Achieves 98.5\% of MBD with 1,000
trajectories; proves pointwise consistency of kernel score estimate.
\textbf{Confirms field has split into model-only vs.\ data-only.}

%---------------------------------------------------------------------
\subsection{Flow Matching for Robot Planning}

\textbf{Flow Matching}~\cite{lipman2023flow}: the foundational FM framework.
Simulation-free CNF training via conditional probability paths. Outperforms
diffusion paths on ImageNet.

\textbf{FlowMP}~\cite{flowmp}: second-order FM for robot motion.

\textbf{SafeFlow}~\cite{safeflow}: flow matching barrier functions.
Training-free safety at test time.

\textbf{SafeFlowMatcher}~\cite{yang2025safeflowmatcher}: prediction-correction
integrator with CBF-QP.

\textbf{UniConFlow}~\cite{yang2025uniconflow}: unified equality/inequality
constraint enforcement in FM.

%---------------------------------------------------------------------
\subsection{Guided Generation via Optimal Control}

\textbf{OC-Flow}~\cite{ocflow}: training-free guided FM via optimal
control (PMP-based co-state equations). Applied to image and molecule generation.
Open-source. Our model-based score is a special case of OC guidance derived
from dynamics rather than a terminal reward.

\textbf{SOCM}~\cite{domingoenrich2024socm}: IDO technique for SOC via
least-squares matching; path-wise reparameterization trick; lower error than
existing IDO methods.

%---------------------------------------------------------------------
\subsection{Safety in Generative Planning}

\textbf{SafeDiffuser}~\cite{xiao2025safediffuser}: CBFs integrated into
denoising; evaluated on maze, legged locomotion, 3D manipulation.

\textbf{S$^2$Diff}~\cite{s2diff}: Lyapunov-guided diffusion model for
safe and stable control. Discovers theoretical link between diffusion sampling
and Almost Lyapunov theory. Eliminates reliance on complex solvers and
control-affine structures. NeurIPS 2025.

%---------------------------------------------------------------------
\subsection{Hybrid and Compositional Approaches}

\textbf{Classifier Guidance}~\cite{dhariwal2021diffusion}: the original
score-level product-of-experts for conditional generation.
$\nabla\log p(x_t|y) = \nabla\log p(x_t) + \nabla\log p(y|x_t).$
Conceptual ancestor of our approach.

\textbf{CFGRL}~\cite{cfgrl}: classifier-free guidance as policy
improvement operator. $\pi(a|s)\propto\hat{\pi}(a|s)\cdot f(A(s,a))$.
Same product-of-experts math in action space; confirms generality of the
composition principle.

%%%%%%%%%%%%%%%%%%%%%%%%%%%%%%%%%%%%%%%%%%%%%%%%%%%%%%%%%%%%%%%%%%%%%%%%%%%%%%%%
\section{Open Questions and Risks}
\label{sec:openquestions}

\textbf{Fairly confident:}
The product-of-experts score decomposition is well-defined;
the LQ-Riccati connection is verifiable analytically;
the $\lambda$-sweep experiment will produce interpretable results;
MBD's Monte Carlo score estimation works (NeurIPS 2024).

\textbf{Uncertain:}
Whether the hybrid ODE converges to the correct product distribution in the
nonlinear case (may require hard-to-verify regularity conditions);
whether Monte Carlo score estimation is accurate enough in high dimensions;
whether the product-of-experts assumption is appropriate (calibration between
data and model distributions may be off);
sensitivity to $\lambda$ (may require per-task tuning).

\textbf{Fallback plan:}
If general theory is intractable, frame around (1) the LQ-Riccati connection,
(2) empirical $\lambda$-tradeoff characterization, and (3) failure mode analysis.
This remains a complete, publishable-quality course project.

%%%%%%%%%%%%%%%%%%%%%%%%%%%%%%%%%%%%%%%%%%%%%%%%%%%%%%%%%%%%%%%%%%%%%%%%%%%%%%%%

\bibliographystyle{IEEEtran}
\bibliography{IEEEabrv,references}

\end{document}
